\subsection{Decimocuarta iteración de la Fase 3: Pruebas Finales}
\label{subsec:VerificacionF3}

Como en las fases anteriores, el cierre de la fase se certifica con una
iteración de validación que recorre todos los cambios incorporados, desde la
carga del cartucho hasta la presentación en pantalla con el GWI real del

\subsection{Batería con las pruebas efectuadas y su resultado}
\label{sec:BateriaPruebasFase3}

% Configuración para que las columnas X justifiquen a la izquierda en lugar de centrado/justificado completo
\newcolumntype{L}{>{\raggedright\arraybackslash}X} 

\begin{small} 
	\begin{tabularx}{\textwidth}{c >{\hsize=0.4\hsize}L >{\hsize=1.6\hsize}L c}
		\toprule
		\textbf{Nº} & \textbf{Grupo} & \textbf{Prueba} & \textbf{Resultado} \\ 
		\midrule
		\endfirsthead % Cabecera de la primera página
		
		\toprule
		\textbf{Nº} & \textbf{Grupo} & \textbf{Prueba} & \textbf{Resultado} \\ 
		\midrule
		\endhead % Cabecera que se repite en las páginas siguientes
		
		\midrule
		\multicolumn{4}{r}{\small Continúa en la página siguiente...} \\
		\endfoot % Pie de página en páginas intermedias
		
		\bottomrule
		\endlastfoot % Cierre definitivo al final de toda la tabla


	1 & Jugador & El jugador queda visible y presentado en su posición de
	inicio al primer \textit{iniciar}, sin necesidad de pulsar nada. & OK \\
	\hline
	2 & Jugador & Los botones mueven al jugador por sus enlaces de
	dirección; los enlaces inversos se mantienen automáticamente al
	marcarlos y al quitarlos. & OK \\
	\hline
	3 & Jugador & La verificación al guardar bloquea un cartucho con un
	jugador sin ningún enlace de dirección a otro jugador. & OK \\
	\hline
	4 & Rutas & Los segmentos de juego nacen por azar, avanzan por sus
	enlaces un segmento por tick y mueren al llegar al final de su
	concatenación. & OK \\
	\hline
	5 & Azar & Con la misma semilla, dos partidas son idénticas: mismas
	rutas naciendo en los mismos ticks y mismos desenlaces. & OK \\
	\hline
	6 & Colisiones & Una colisión solo se decide entre el jugador y un
	segmento visible que reporta colisión; el desenlace salvado enciende su
	cadena y el mortal la suya. & OK \\
	\hline
	7 & Congelación & Mientras corre una cadena de desenlace el juego queda
	congelado: no nacen rutas y el jugador no responde; al morir el último
	eslabón el juego se reanuda. & OK \\
	\hline
	8 & Contadores & La familia de vida se enciende desde el primero y la de
	muerte acumulativamente desde el último, en el orden definido por su
	cadena y no por la numeración del diseñador. & OK \\
	\hline
	9 & Contadores & Los contadores de dígitos presentan su valor con
	segmentos de siete segmentos y lo actualizan en el mismo tick del
	cambio. & OK \\
	\hline
	10 & Fin de partida & La partida termina al agotarse un contador, al
	encenderse todos los iconos de muerte o al perderse todas las vidas,
	según el mecanismo declarado en el cartucho. & OK \\
	\hline
	11 & Adicionales & Un segmento \textit{adicional} sin siguiente queda
	fijo encendido (helicóptero); con siguiente avanza en su propio bucle, y
	un \textit{inicioAdicional} nace encendido y alterna solo (aspas), sin
	que nada los interrumpa. & OK \\
	\hline
	12 & Secuencia & La Secuencia lanza sus segmentos en orden, uno por
	tick, y se relanza desde el primero tras esperar los ticks declarados
	(ciclo completo de $n+m$ ticks). & OK \\
	\hline
	13 & Secuencia & Las colisiones y la congelación por desenlace salvado
	NO interrumpen la Secuencia; solo el nacimiento de un ciclo de muerte
	la interrumpe. & OK \\
	\hline
	14 & Secuencia & Al nacer el ciclo de muerte, el segmento activo de la
	Secuencia desaparece, la animación se reinicia por completo y espera
	los ticks necesarios para relanzarse desde su primer segmento. & OK \\
	\hline
	15 & Secuencia & La cadena de muerte atraviesa un segmento encendido
	por la Secuencia (las aletas compartidas): lo toma como eslabón, lo
	apaga y sigue, produciendo la muerte. & OK \\
	\hline
	16 & Acompañantes & Un secundario «Mostrado junto a» se enciende y
	apaga siempre con su principal, incluidas las cadenas de desenlace, y
	puede presentarse solo en la animación de fondo. & OK \\
	\hline
	17 & Acompañantes & La verificación al guardar bloquea un secundario
	con cualquier enlace: con siguiente, apuntado por otro, con colisión, de
	tipo distinto a juego o con dos principales. & OK \\
	\hline
	18 & Depuración & La barra de estado con \texttt{GW\_DEBUG} muestra el
	estado de la Secuencia (contador pendiente, segmento en curso y si
	anima o está suspendida). & OK \\
	\hline
	19 & Integración & Partida completa con el GWI real del Parachute: tres
	ciclos de muerte que atraviesan las aletas encendidas por la
	animación, tres iconos de muerte y partida terminada por fin de
	muerte. & OK \\
	\hline
	
	\end{tabularx}
	\nopagebreak
	\captionof{table}{Batería de pruebas de validación de la Fase 3.}
	\label{tab:PruebasFase3}
\end{small}